\section{Introduction}
\label{sec:intro}

Deployed language-model agents are specialized systems
\citep{yao2023react,schick2023toolformer,jimenez2024swebench,liu2024agentbench} whose cost scales
with each token. The model that serves such a workload is usually a reference model scaled down by pruning,
quantization, or distillation, and choosing the method and its strength currently requires
measuring each candidate. A single compression ratio does not replace that search, because
compression does not reduce capabilities uniformly: across twelve public models the same pruning density leaves QA the
least damaged capability in some series and the most damaged in others
(Fig.~\ref{fig:heterogeneity}), and the ordering changes with size and family, so capability must
enter the prediction.

We ask how the capability-conditioned loss after compression can be predicted from the initial
model state and the configuration of the method. The endpoint is the conditional loss of each
capability on a fixed probe set relative to a reference measured beforehand
(Eq.~\ref{eq:interface}); the admitted inputs are the initial state (parameter count, pretraining
tokens, reference capability losses) and the intervention setting, and no measurement of the
compressed target is used. Relations are developed on a controlled panel of Pythia checkpoints and
tested with coefficients and predictions frozen in advance.

The study proceeds in one sequence, from the prediction problem and its admitted inputs, through
controlled measurement and explanation, to frozen tests and a selection decision. Each relation,
verdict and recommendation keeps the status assigned before its outcome was measured: development,
frozen prediction, or post-hoc recommendation.

\paragraph{Contributions.} Three findings, and what they enable, make up the contribution.
\begin{itemize}
\item \textbf{Capability responses differ in shape, and the shape decides what a law may share.}
Math and code follow one shared pruning shape with capability-specific scales, while the
question-answering response changes sign across density and cannot be a scaled copy of that curve;
a coefficient-free second-order account of logit displacement explains the measured damage inside a
stated empirical range.
\item \textbf{Predicting new configurations of measured models and transferring to new model
states need different inputs and evidence.} Initial-state inputs carry math and code to unseen
densities and group sizes of measured states; on new states a source-free configuration curve is as
good or better, and a rectangle design measures the budget-by-reuse interaction.
\item \textbf{The data requirement of distillation is defined by training exposure, student and
evaluation distribution together.} Under one fixed protocol the learning rate sets the size of the
damage, the cost of reusing a small pool grows with the reuse ratio and the student size and
reproduces across three question-answering distributions, and the gain appears on one.
\item \textbf{What the findings enable.} Method-specific relations with their tested ranges
(Table~\ref{tab:main-prediction-v2}), a selection rule whose frozen maps on four fresh Pythia
states match the measured oracle within 0.00001, 0.004 and 0.14 nats for math, code and QA, and a
released suite of coefficients, registered predictions, measurements and code.
\end{itemize}
